\documentclass[../main.tex]{subfiles}
\begin{document}
%==========================================TIÊU ĐỀ TẬP========================================
\begin{table}[h]
    \begin{adjustwidth}{-1cm}{}
        \begin{tabular}{>{\centering\arraybackslash}p{6cm}|>{\raggedright\arraybackslash}p{12.5cm}}
            \multirow{5}{6cm}
            % Cột bên trái: ÉPISODE và số tập
            {\rainbowcolor
                {\\[-40pt]
                    \flaregothic\fontsize{20pt}{20pt}\selectfont \centering
                    \color{eptype}JOUR\color{black}\\[12pt]
                    \hbrk\fontsize{72pt}{72pt}\selectfont
                    \color{epnum}4
                    \\[18pt]
                    % \flaregothic\fontsize{14pt}{14pt}\selectfont
                    % \color{epattr}BIENVENUE, \uline{\color{Banpire} Mel} !
                    % \flaregothic\fontsize{18pt}{18pt}\selectfont
                    % \color{epattr}ÉPISODE PERDU\\
                    \unrainbowcolor}
                \color{black}}
             &
            % Dòng thứ nhất: tên tập tiếng Nhật
            {\fotskip\fontsize{18pt}{18pt}\selectfont おもいでいっぱい！
            }                                                   \\[6pt]
            % Dòng thứ hai: tên tập tiếng Anh
             & {\cafeteria\fontsize{18pt}{18pt}\selectfont Moving Day with VTubers}                              \\[4pt]
            % Dòng thứ ba: tên tập tiếng Việt
             & {\vnmsans\fontsize{18pt}{18pt}\selectfont Ôn lại kỷ niệm}                                         \\[8pt]
            % Dòng thứ tư: Nhân vật xuất hiện
             & {\sen{\fontsize{14pt}{14pt}\selectfont Track used: Relieve in Paradise - ThirstyHammer \& Lexay}}
            \\[6pt]
            % Dòng cuối cùng: Thời gian diễn ra của tập (thường sẽ là ngày phát hành)
             & {\continuum\fontsize{16pt}{16pt}\selectfont Ngày phát hành: 12/06/2021}                           \\[18pt]
        \end{tabular}
    \end{adjustwidth}
\end{table}
%=========================================TRUYỆN CHÍNH========================================
\color{black}

\pov{Hikawa Kuon}

Tôi đứng trước khu văn phòng Nijisanji cũ, nơi mà chỉ mới hôm qua còn là trụ sở làm việc, giờ thì đã trở thành một đống gạch vụn đúng nghĩa. Một phần bức tường bên trái còn đang bốc khói, trông chẳng khác gì một cảnh hậu chiến. Tôi khoanh tay, nhìn nó thêm vài giây, rồi thở dài: ``\textit{Làm mình nhớ tới hồi mà Ayame-chan hay Robo-PON phá sập cả cái văn phòng bên \hll\ hai năm về trước\footnote{Xem lại lần lượt ở các tập {\flaregothic ÉPISODE 3}, quyển số Không và {\flaregothic ÉPISODE 22}, quyển số Hai.}\dots}''

Ừ thì, ở thế giới bên kia, mấy cái chuyện nổ văn phòng, vỡ kính, cháy bảng mạch cũng không còn lạ nữa. Nhưng lần này thì hơi quá thật, cho dù đây không phải văn phòng của chúng tôi. Và hậu quả thì cũng chẳng ai ngờ được.

Sau cú nổ ``kinh thiên động địa'' đó do lỗi của  Rena-chan, cả trụ sở Nijisanji buộc phải chuyển sang một nơi mới cách đó không xa lắm. Vẫn là một khu tòa nhà văn phòng, chỉ là\dots\ hơi tạm bợ hơn bình thường.

Bên trong, vừa bước vào cái hành lang chưa sơn xong, tôi đã nghe tiếng trò chuyện ồn ào từ mấy người quen ở quán ăn ngày đầu.

\begin{figure}[H]
    \centering
    \begin{subfigure}[t]{0.45\textwidth}
        \centering
        \includegraphics[width=8cm]{../../../../Image/Book?/nj-4a1.png}
        \caption{Vị trí văn phòng Nijisanji cũ.}
    \end{subfigure}
    \quad
    \begin{subfigure}[t]{0.45\textwidth}
        \centering
        \includegraphics[width=8cm]{../../../../Image/Book?/nj-4a2.png}
        \caption{Vị trí văn phòng Nijisanji mới, với Dola, Claire và Ryuushen ăn mừng.}
    \end{subfigure}
\end{figure}

``Đến thời điểm hiện tại thì\dots'' Claire cười nhẹ nhàng, như thể đang đọc lời cầu nguyện ở nhà thờ.

``NIJISANJI đã chuyển trụ sở sang văn phòng mới!'' Dola reo lên.

``Và nó cách chỗ cũ không xa!'' Ryuushen thêm vào.

Cả ba người đồng thanh, cùng giơ tay ăn mừng: ``\textbf{Yay!!}''

Tôi nghe mà chỉ muốn gục mặt xuống sàn. Mệt rã rời, tôi vẫn cố cất giọng cho đủ phép lịch sự: ``\dots Yay dủng cái nỗi gì. Vác đồ mệt gần chết luôn.''

Cả ngày hôm nay, tôi phải dọn đống đồ từ cái văn phòng vừa bị thổi bay hôm qua sang đây. Tất cả các hộp đều là của người khác: máy quay, dụng cụ stream, đạo cụ, mấy bộ trang phục lạ lùng\dots\ Đừng hỏi tại sao tôi phải làm một mình - vì tôi cũng không biết. Kaigou-chan và Suzuki thì vẫn đang ngủ ở nhà, Game-san thì vẫn đang sạc, còn bản thân tôi chỉ đang đi hóng mát thì bị bắt làm việc này.

\begin{wrapfigure}[29]{r}{4.5cm}
    \includegraphics[width=4.5cm]{../../../../Image/Book?/claire_model.png}
\end{wrapfigure}

Để tôi giới thiệu ba con người này thuộc nhánh SEEDs một chút nhỉ?

\chrint

Người đầu tiên là Claire (クレア, Cư-rê-a), hay chính xác hơn là Chị sơ Claire (シスター・クレア), là một người - chắc bạn cũng đoán được rồi - làm ở trong các nhà thờ, hay tham gia các hoạt động thiện nguyện. Vì là một bà sơ nên rất quý trọng Chúa, và là một người có thể gọi là ``ngây thơ'' và trong sạch. Nói thật thì lúc mới gặp tôi cứ tưởng cô là người ngoài ngành - nhưng mà cái cách cô gật đầu mỗi khi tôi càu nhàu lại khiến tôi thấy dễ chịu.

Mặc dù là một người theo đạo, cô nàng lại có máu hài hước rất mạnh mẽ, đôi lúc còn ô uế hình ảnh hướng thiện Chúa của mình mà không hề hay biết. Tôi nghe nói cô ấy còn làm những điều gì đó mà khiến không chỉ tôi và mọi người, mà thậm chí đến vị Chúa trên cao còn phải muốn thở dài ngán ngẩm.

Claire-san rất thích hát, đồng thời có một quãng giọng rất lớn, tuy thế lại thích hát những bài ở nốt cao hơn là nốt trầm.

Cô ấy được nhận xét là một người ngăn nắp và cẩn thận, mặc dù bản thân cô không nghĩ mình là một người như vậy, thậm chí đôi lúc còn tự phụ tới mức khó tin.

Dù là một bà sơ, cô nàng cũng co những khoảng thời gian được sống và làm việc như một người bình thường. Cô nàng có thể chơi game buổi đêm, trò chuyện vào lúc người ta đã ngủ, và cũng có hứng thú đôi chút về hóa trang.

Claire-san có mái tóc dài màu nâu nhạt (hoặc vàng, tùy theo cách bạn nhìn nhận), với mái thưa rẽ ngôi và đôi mắt nâu. Cô ấy mặc bộ trang phục thường thấy ở các nữ sơ ở nhà thờ: áo choàng trắng, váy dài màu đen cao cổ, có tay áo dài và khóa kéo đằng sau. Cô đeo găng tay trắng, quần tất đen dài đến đầu gối có dây nịt đi và giày lười màu nâu.

\begin{wrapfigure}[29]{l}{11cm}
    \includegraphics[width=11cm]{../../../../Image/Book?/dola_model.png}
\end{wrapfigure}

Người tiếp theo là Dola (ドーラ, Đô-ra), một con rồng lửa hơn 350 tuổi mang hình dạng giống như một cô gái. Vốn dĩ cô nàng có nhiệm vụ canh giữ một hang động kho báu ở thế giới khác, nhưng một thứ gì đó khiến sẽ cô nàng bị xao nhãng, và thế là cô tới đây để gia nhập Nijisanji luôn.

Cô nàng này được coi là ``cầm đầu'' của cả nhánh SEEDs, khi thậm chí bản thân cô là một trong những người sáng lập ra server (tôi cá là Reddit) và mang cho mình chức quản trị viên ở đó. Dù vậy, cách nói chuyện hơi ngọng nghịu của cô khiên cho cô trở thành một trong những người dễ tính nhất trong nhánh.

Dola-chan là một trong những thành viên kỳ cựu nhất của Nijisanji, nổi bật với tính cách năng động, thân thiện và khả năng lãnh đạo tự nhiên. Cô thường được xem là người gắn kết các thành viên trong nhóm, luôn chủ động tổ chức các hoạt động chung và hỗ trợ mọi người khi cần thiết. Cô rất nhiệt tình tham gia các dự án hợp tác, từ các buổi stream nhóm cho tới những sự kiện lớn của công ty.

Ngoài khả năng ca hát và nhảy múa, cô còn có kỹ năng chơi game khá tốt, đặc biệt là các trò chơi hành động và sinh tồn. Cô thường xuyên xuất hiện trong các buổi stream với tinh thần vui vẻ, hài hước, tạo không khí thoải mái cho cả người xem lẫn các thành viên khác.

Cô ấy cũng nổi tiếng với sự kiên nhẫn và khả năng xử lý tình huống bất ngờ. Trong những lúc nhóm gặp khó khăn hoặc có sự cố, cô luôn là người đứng ra giải quyết và động viên mọi người tiếp tục cố gắng. Nhờ vậy, Dola được nhiều người yêu quý và kính trọng, không chỉ trong nhánh SEEDs mà còn trong toàn bộ cộng đồng Nijisanji.

Dola-chan có bộ trang phục lộ liễu nhất mà tôi từng thấy. Cô chỉ đeo một vòng cổ đen và một dây đeo ở đùi phải. Phần lớn thứ che đi những bộ phận nhạy cảm nhất đều là vảy rồng lửa mang ba màu sắc cam, vàng và đỏ, cùng đôi guốc (?) cao gót và chân trái che gần hết. Cả hai cánh tay của cô cũng là vảy rồng, và thứ duy nhất trông như da người là chân phải, phần khe ngực và bụng. Cô ấy còn có một cái đuôi dài phía sau. Về gương mặt, Dola-chan có mái tóc đỏ dài để xõa, măt đỏ, sừng rồng - hao hao giống như Coco-chan ngày nào.

\begin{wrapfigure}[29]{r}{8cm}
    \includegraphics[width=8cm]{../../../../Image/Book?/ryuushen_model.jpg}
\end{wrapfigure}

Người cuối cùng là Ryuushen (\ruby{緑}{りゅう}\ruby{仙}{せん}, \ruby{Ri-ư}{Lục}-\ruby{sen}{Tiên}), tên thật là Sengawa Midori (\ruby{仙}{せん}\ruby{河}{がわ}\ruby{緑}{みどり}, \ruby{Sen}{Tiên}-\ruby{ga-oa}{Hà}\ \ruby{Mi-đô-ri}{Lục}). Em ấy\dots\ à không, đằng ấy (tôi vẫn chưa quen cách xưng hô với người không công khai giới tính) là người có phong cách rất truyền thống, mà tôi không rõ là người ấy đang mặc bộ trang phục kiểu Nhật Bản hay Trung Quốc nữa. Một người mà tôi có thể nói vui là ``xăng pha nhớt'' - tức là mập mờ giới tính.

Ryuushen chơi rất thân với Dola-chan và Claire-san - mà đúng hơn, là cả ba người họ chơi thân với nhau - tới mức lập ra một nhóm riêng, và còn sáng tác ra bài hát dành cho nhóm.

Đằng ấy là một người có tính cách hơi lập dị, khiến cho việc kết bạn ở trường hồi trung học khá khó khăn, và đó là lý do vì sao mục tiêu đầu tiên của người ấy là kết bạn 100 người. Dù là một người không tiết lộ giới tính, Ryuushen rất tự tin về ngoại hình và khả năng ca hát của mình.

Đằng ấy có sở thích mãnh liệt với các bài hát và anime vào thập niên những năm 80, và thường gắn bó chặt chẽ hơn với những thành viên lớn tuổi qua những sở thích chung này.

Cùng với Chaika-san, đằng ấy là một người trong nhóm ``kháng chiến Nijisanji'', nhưng chỉ là trên giấy tờ, và thường không xuất hiện thường xuyên trong nhóm đó. Ryuushen cũng là một người khá kém về khoản chơi game, nhưng có thể thực hiện tốt các buổi phát sóng về trò chuyện.

Đằng ấy cũng là một tay vẽ giỏi. Công nhận là như thế thật.

Ryuushen có đôi mắt tím, mái tóc xanh lục dài được tết lệch bên trái, trên đỉnh có chỏm tóc ngố. Đằng ấy - xin lỗi nếu nghe không thuận tai - mặc áo cụt tay màu đen bên trong, vòng cổ màu nâu, găng tay đen bó sát, áo khoác trắng mặc ngoài (thường để trễ vai) với tay áo dài tới độ che cả bàn tay. Phía trước uần đen đeo khố in họa tiết màu xán lớn, đi giày bệt màu đen và không đi tất. Ở cả hai tai người này xỏ nhiều khuyên, trong đó có khuyên tua rùa.

\chrint

Cả ba đang đứng trong một căn phòng trống hoác, bên trong mới chỉ lắp đặt những tiện ích cơ bản, bên cạnh là một đống hộp mà tôi vừa còng lưng khuân về.

\img{nj-4b}

Claire-san ngó nghiêng khắp nơi, nói một cách\dots khac lạc quan: ``Bọn mình đã tới đây cho buổi ghi hình đầu tiên trong văn phòng mới này!''

``\dots Tại sao mọi chuyện lại trở nên như thế này!?'' Ryuushen-san thì hoảng hốt, mắt đảo quanh.

Tôi thì chỉ còn đủ sức ngồi bệt xuống đất, vừa thở vừa lau mồ hôi. ``Anh vừa mới chuyển mấy cái hộp này từ văn phòng cũ bị nổ hôm qua sang đấy.''

Dola-chan nhìn tôi chằm chằm: ``Khoan đã, anh làm hết chỗ đó một mình sao?''

``\dots Ừ.'' Tôi ngả đầu ra sau. ``Kaigou-chan với Suzu thì vẫn đang say giấc nồng, còn Game thì anh cũng không gọi được vì quên không sạc pin\dots\ Vãi chưởng mệt quá!''

Người mê Gấu trúc ngờ vực: ``Anh nói `nổ' là sao?''

Tôi bực mình gác tay lên gối, rồi thở dài kể lại: ``Mấy đứa không biết à? Chỉ vì tên ngốc nào đó lấy Chaika-san làm con tin để đòi có được cái cài tóc `siu đáng iu\~{}'. Cuối cùng là vụng tay vụng chân một cái, \textbf{BÙM!} Thế là bên công ty phải dọn luôn sang đây.''

``Tội cho Kuon-san thật\dots'' cô Nữ sơ thì thở dài.

``Eh, không cần phải lấy làm tội cho anh đâu,'' tôi khoát tay. ``Quen rồi ấy mà. Giờ tí còn phải mở mấy cái hộp này để sắp xếp lại này.''

``Ít ra thì nếu anh cần giúp chuyển hết số hộp đó thì cũng phải nói với tụi này một tiếng chứ,'' Ryuushen-san trách, ``Bên công ty không thông báo gì à?''

``Chắc là vì họ nghĩ bọn mình sẽ không tới nếu nói cho bọn mình đấy\dots'' Claire-san nói nhỏ.

``Kiểu gì tụi mình cũng sẽ tới thôi!'' đằng ấy phản bác.

Tôi lắc đầu: ``Anh không nghĩ đó là lý do đâu. Có lẽ người ta thấy ai giúp được là cứ thế hốt vào thôi.''

Người tóc xanh chỉ ``hừm'' một tiếng, rồi thôi không tranh luận nữa: ``Mà thôi. Để đấy để bọn này phụ giúp anh vậy.''

``Cứ việc.'' Tôi ngồi thẳng dậy, tay chỉ về đống hộp. ``Có gì khó thì cứ nhờ anh.''

Claire-san nhẹ nhàng cười trừ: ``Anh không cần phải khệ nệ trước bọn em đâu\dots''

Tôi nhún vai, gật đầu, và ngáp một cái rõ dài. Mắt vẫn còn nheo lại, nhưng lòng lại thấy nhẹ đi phần nào. Có người giúp, vẫn hơn là làm một mình giữa cái tòa nhà trống rỗng này.

\ct

``Trước tiên thì, tại sao chúng ta không đưa hết tất cả mọi thứ ra ngoài trước nhỉ?'' Dola lên tiếng, tay chống hông, mắt nhìn quanh căn phòng rồi lại nhìn về phía tôi với Ryuushen và Claire.

``Ý hay đấy,'' tôi gật đầu, ngồi bệt trước chiếc quạt điện mà mình vừa moi ra được từ một trong mấy cái hộp. Gió phả vào mặt thật dễ chịu, nhưng mồ hôi vẫn cứ túa ra như tắm.

``Vậy thì mình sẽ mở hộp này nha\~{}'' cô Bà sơ lên tiếng, nở nụ cười nhẹ nhàng như thường lệ. Người tóc xanh thì không nói gì, chỉ gật đầu rồi ngồi xổm xuống trước một chiếc hộp khác, bắt đầu sục sạo bên trong.

``Trong này có cái gì nào\dots'' đằng ấy lẩm bẩm, rồi đột nhiên dừng lại một chút. Tôi thấy vai Ryuushen hơi khựng lại. ``Hmm?''

\begin{wrapfigure}{r}{6cm}
    \includegraphics[width=8cm]{../../../../Image/Book?nj-4_1.png}
\end{wrapfigure}

Tôi nghiêng đầu nhìn theo. Ryuushen từ từ lôi ra một tấm gì đó, hóa ra là một tấm poster. Đằng ấy giơ cao lên, mặt sáng bừng lên vì phấn khích.

``Nhìn này, nhìn này! Đây là tấm poster hồi tụi mình đi lưu diễn ở VARK đó!'' đằng ấy reo lên.

Tôi nheo mắt, hơi nghiêng đầu để nhìn kỹ hơn. Trên tấm poster là hình ba người họ: Ryuushen, Dola và Claire. Tên nhóm ``cresc.'' hiện rõ phía góc trên, bên cạnh dòng chữ về buổi concert đầu tiên trong VARK.

``Giờ anh mới biết là mấy đứa cũng có concert đấy,'' tôi thốt ra thành thật, vẫn còn ngạc nhiên. ``Anh cứ tưởng mấy đứa chuyên stream game thôi chứ?''

Ryuushen đặt tấm poster dựa vào thùng giấy, nhún vai. ``Thỉnh thoảng bọn em cũng ra mát một vài bài hát nữa. Thậm chí còn làm cả buổi live show hẳn hơi, như thế này chẳng hạn.''

``Ra là thế\dots'' tôi gật gù. ``Về cơ bản thì cũng học theo bên bọn anh đấy thôi.''

``Ối chà,'' nàng Rồng lửa đột ngột phá lên, ``làm mình nhớ hồi đó ghê á. Buổi concert đó vui cực luôn!''

``Quả nhiên!'' cô Nữ sơ cười khúc khích. ``Ước gì bọn mình có thêm một buổi giống vậy nữa, với ba chúng mình nhỉ.''

``Đúng ha\~{}'' Dola đồng tình rồi quay sang tôi, gọi bằng cái biệt danh mà mấy người bên đây vẫn thường dùng, mà có lẽ chỉ có một người cùng loài bên \hll\ hay gọi. ``Mà, Kuon-paisen?''

``Sao?''

``Bên bọn anh chắc là cũng có nhiều buổi biểu diễn như thế nhỉ?''

``Ờm, đương nhiên.'' tôi gật đầu, nhún vai như đó là lẽ thường. ``Bên \hll\ là lò đào tạo idol mà. Thế nên bọn anh mới tổ chức nhiều thế chứ. Có cả online lẫn offline luôn.''

``Bọn anh cũng có online sao?'' Ryuushen hỏi, hơi nghiêng đầu.

``Ừ. Cái buổi live show \textit{Bloom,} hồi COVID năm ngoái là một ví dụ.'' Tôi liếc qua cái quạt, vẫn còn xoay đều đều trước mặt mình, cảm thấy cái tên đó vừa bật ra lại khiến tôi nhớ tới hàng đống kỷ niệm và cả hàng tấn công việc đằng sau nó. ``\hl{Cũng là buổi live mà mình không tham dự được do ăn Tết.}''

``Vậy\dots'' Claire-san nói với vẻ chần chừ nhưng có phần hứng khởi, ``sao anh không thử xem một trong những buổi live show từ bọn em nhỉ?''

``Ờm\dots'' Tôi ngập ngừng, tay gãi nhẹ sau đầu. ``Cái này thì\dots\ hơi khó. Bận bù đầu thế này còn đâu ra thời gian mà xem nữa\dots''

``Vậy à\~{}\dots'' Dola-chan nheo mắt, giọng hơi kéo dài như kiểu không mấy tin tưởng.

Tôi chỉ biết cười trừ. Sự thật là tôi có thể xem nếu muốn. Nhưng mà khi tôi đã cùng với mọi người thực hiện tới cả trăm buổi concert bên \hll\ rồi, việc ngồi xem người khác diễn - dù là bên Nijisanji - tự nhiên lại thấy lạ. Không hẳn là ngại, cũng không phải coi thường. Chỉ là\dots\ nó không giống với không khí mình quen thuộc. Nó như thể tôi đang nhìn một vũ trụ song song vậy.

\ct

Ba người bọn họ tiếp tục moi đồ ra khỏi mấy cái hộp, vừa sắp xếp, vừa trò chuyện rôm rả. Tôi thì tranh thủ ngồi duỗi chân một chút. Dù nói là ``phụ giúp'' nhưng từ nãy giờ tôi là người bê nặng nhiều nhất, còn giờ thì ai lo hộp nấy.

``Tiếp theo sẽ là\dots'' Ryuushen lầm bầm, rồi cúi xuống, khua khua trong chiếc hộp người đó đang xử lý. ``A, đây rồi.''

Đằng ấy kéo ra ba cái standee lớn 0 kiểu khung treo hình nhân vật - in hình ba người họ: Claire, Dola và chính đằng ấy. Hình như là loại đã ra mắt cách đây không lâu.

\begin{wrapfigure}{l}{8cm}
    \includegraphics[width=8cm]{../../../../Image/Book?/nj-4_2.png}
\end{wrapfigure}

``Có vẻ như là một trong những mặt hàng mới được bán gần đây,'' Ryuushen vừa nói vừa dựng thử lên.

``Ồ\~{}'' Dola-chan reo nhẹ. ``Mình vui ghê khi có những sản phẩm merch mà có mặt cả ba đứa tụi mình trên đó!''

Claire-san cũng cười hiền khi nhìn lại thứ này. ``Bên bọn anh chắc là cũng có những thứ thế này, nhỉ, Kuon-san?''

``Ờm, tất nhiên rồi,'' tôi gật đầu, quẹt mồ hôi. ``Thậm chí nhiều hơn thế nữa là đằng khác. Không chỉ mỗi standee đâu, còn có cả CD, đĩa đơn, bài hát, ảnh, sticker, vân vân mây mây. Nói sao nhỉ\dots\ mấy đứa biết đấy, idol mà.''

``Cái đó thì bọn em hiểu,'' Ryushen gật đầu, ``nhưng mà, còn gì đặc biệt hơn không?''

Tôi đưa tay gãi đầu. ``Hmm\dots\ Gói giọng (voice pack), móc chìa khóa, huy hiệu, huy chương, mấy thứ lặt vặt ấy. Và không chỉ dành riêng cho các buổi live show đâu. Mỗi thành viên ở bên anh thỉnh thoảng cũng có merch riêng. Còn mấy dịp như sinh nhật hay lễ kỷ niệm ra mắt thì khỏi nói, đặc biệt nhiều luôn. Thậm chí còn có theo nhóm.''

``Cái này thì mới đó nha\~{}'' nàng Rồng lửa vỗ tay, mắt sáng lên.

Tôi chỉ nhún vai. ``Chả phải đâu. Có từ hồi\dots\ trời, hồi Tokino Sora-chan còn solo một mình cơ, lúc anh còn chưa vào ấy. Lâu lắm rồi.''

Ryushen bật cười. ``Cái cách anh nói kiểu như già hơn bọn em mười tuổi ấy.''

``Thì\dots'' tôi rướn vai, ``anh là senpai mà. Anh cũng là holoStaff chủ chốt bên đó luôn.''

Mọi người phá lên cười. Claire còn dịu dàng thêm vào một câu: ``Thế nên mới đáng tin cậy thế này, nhỉ.''

Tôi cười gượng. Tôi chưa từng nghĩ sẽ có ngày mình được mấy đứa bên Nijisanji gọi là ``đáng tin cậy'', nhất là sau cái vụ nổ tung văn phòng hôm trước\dots\ Thật là.

\ct

Công cuộc dọn dẹp vẫn tiếp tục. Ryushen lại thọc tay vào hộp, lần này là một chiếc khác đặt ở góc phòng.

``Còn gì nữa đây\dots \textbf{À!}'' đằng ấy reo lên rồi moi ra một chồng sách nhỏ. Nhìn bìa là tôi đoán được ngay - fanbook.

\img{nj-4_3}

``Đây là những cuốn fanbook tụi mình làm hồi lễ hội đầu năm nay á,'' Ryushen vừa nói vừa giơ một quyển lên cho tôi xem.

Tôi cầm lấy một cuốn, lật thử vài trang. Màu sắc rực rỡ, ảnh các thành viên bên này, chữ ký, lời nhắn, bình chọn\dots\ khá giống với cái format mà bên \hll\ cũng từng làm. Nhưng nét trình bày có phần riêng biệt, mang cảm giác ``Niji'' rõ ràng là tươi sáng, tự do, bộc trực hơn.

``Hình như bên bọn anh cũng có vài cuốn kiểu vậy\dots'' tôi vừa nói vừa gật gù, tay vẫn giở sách.

``Ồ!'' Claire-san ngạc nhiên. ``Có được nhiều không ạ?''

``\dots Không được nhiều lắm đâu, anh nghĩ thế.'' Tôi đưa lại quyển sách, thở nhẹ. ``Mỗi lần làm fanbook là một lần cực hình. Quản lý thì giục bài, vẽ thì phải duyệt ba lần, sửa lời nhắn thì sửa đi sửa lại cả chục lượt\dots''

``Nghe giống tụi em phết,'' Dola-chan cười lớn.

Tôi chỉ biết lắc đầu. ``Nghề idol mà. Cái gì cũng phải hoàn hảo cả.''

``Thế nên mới nể anh đó, Kuon-san.'' Ryushen mỉm cười. ``Em nghe nói là anh có mặt trong gần hết các concert của \hll\ mà.''

Tôi cười khan, trong đầu hiện ra vô số kỷ niệm - từ những buổi tổng duyệt căng thẳng cho tới mấy phút đứng đằng cánh gà, chờ ánh đèn sân khấu hạ xuống, chỉ để rồi tự tôi phải rút về sớm vì còn lo những công việc khác.

``Ờ\dots\ nhưng cũng đâu ít lần anh có được cùng mọi người lên sân khấu đâu\dots''

``Ể!?''

``Thôi, chuyện dài lắm,'' tôi phẩy tay. ``Lo xếp đồ cho xong đã.''

\ct

Đến lượt Claire-san loay hoay mở một cái hộp to đùng ở góc phòng. Tôi vừa kê lại cái quạt cho quay hướng về phía mọi người thì nghe thấy cô nàng thốt lên: ``\textbf{Tôi tìm thấy được thứ gì này!}''

Tôi quay sang thì thấy Claire đang lôi ra một cái nồi khổng lồ, trông cứ như cái thùng rác loại nhỏ. Nhìn mặt cô là biết bất ngờ thế nào rồi.

\img{nj-4_4}

``Đây chẳng phải là chiếc nồi mà bọn mình từng dùng để nấu cà ri sao?'' cô nàng hỏi, mắt vẫn tròn xoe.

Tôi ngó đầu nhìn cái nồi: ``Ừm, có vẻ là vậy đấy. Ở bên anh cũng có một cái kiểu giống thế. Nhưng mà nó đa dụng hơn nhiều.''

Ryushen-san tò mò: ``Thế á?''

``Ừ, bọn anh dùng để nấu mochi, rồi chè, rồi nấu nước dùng nữa. Tiện lắm, khỏi phải mua ba bốn loại nồi khác nhau. Mà lại có thể cho nhiều người cùng nấu một lúc.''

``Nghe có vẻ vui nhỉ!'' Dola cười toe toét.

``Vui thì vui, nhưng mà\dots'' tôi cười nửa miệng. ``Lắm thầy nhiều ma em ạ. Mỗi người nêm nếm một kiểu, thi thoảng lại ra một món\dots\ kỳ dị khó tả. May mà vẫn ăn được.''

Dola-chan bật cười: ``Em rất muốn tập hợp mọi người lại để cùng nhau nấu ăn quá!''

``Mừng cho mấy đứa,'' tôi nhún vai. ``Chỉ mong là không nổ bếp.''

\ct

Ryushen-san lại lục thêm đồ trong cái hộp đang mở dở. ``\textbf{A!}'' đằng ấy reo lên.

Tôi quay đầu lại, thì thấy cậu ta đang giơ lên một cái dây nhảy màu tím. Trông như loại đồ chơi cho thiếu nhi ấy.

\img{nj-4_5}

``Đây là chiếc dây nhảy mà tụi mình chơi hồi stream 3D lần trước này!'' Ryushen-san hào hứng.

``À, bọn anh cũng có cái như vậy,'' tôi gật gù. ``Mà giờ chắc nó nằm xó ở đâu rồi\dots''

Claire-san lúc này cũng khom người xuống hộp, tiếp tục tìm kiếm. Còn tôi thì đưa mắt nhìn quanh, thấy phòng bắt đầu bớt lộn xộn hơn một chút.

Tôi định hỏi tụi nó xem có tìm thấy cái dây nguồn nào không thì cô Nàng sơ chợt kêu lên, như sực nhớ ra điều gì đó: ``Thôi chết! Đáng ra bọn mình phải đi sắp xếp lại đồ đạc mới đúng chứ!''

Tôi bật cười nhẹ, lắc đầu. ``Đó, đó là vấn đề đấy,'' tôi nói, không trách gì cả. Thật ra nhìn mọi người vui vẻ ôn lại quá khứ cũng hay ho phết. ``Anh biết cái cảm giác ôn lại kỷ niệm thời xưa cũ mà. Có mấy lần anh dọn kho đồ trong studio bên mình cũng y như thế.''

Ryushen-san đồng tình: ``Cứ mỗi lần tụi mình dọn dẹp là lại như vậy, nhỉ?''

``Đúng á\~{}'' Dola gật gù. ``Kiểu gì bọn mình cũng sẽ lại chìm vào manga cũ, hoặc là ngồi ngắm album ảnh rồi cười một mình.''

Claire-san mỉm cười, có phần nhẹ nhàng hơn: ``Cái đó thì bọn mình hiểu, nhưng mà\dots\ nên để lại mấy chuyện đó sau khi chúng ta xong việc nhỉ?''

Cả ba người kia đồng thanh: ``Rồi\~{}/OK./Được!''

Cô ấy lúc này mới quay lại nhìn tôi, đang ngồi nghỉ được một lúc. ``Kuon-san, nhờ bọn em phụ dỡ đồ nhé?''

Tôi giơ ngón cái lên, ra hiệu ``like''.

``Rồi, tập trung vào việc nào!''

\begin{center}
    \uline{Ba tiếng sau.}
\end{center}

Tôi cũng không rõ là mấy giờ nữa. Chỉ biết là trời đã ngả sang màu vàng cam từ lúc nào. Ánh sáng chiếu xiên qua cửa sổ lớn của phòng, đổ dài trên sàn nhà đầy hộp bìa và standee lắp dở. Mấy tiếng đồng hồ vừa rồi, tôi cùng ba người bọn họ lôi hết đồ ra khỏi hộp, sắp xếp sơ sơ lại vài thứ. Cũng tạm gọi là có hình có dạng.

Tôi ngồi phịch xuống một chiếc ghế bành còn nguyên bọc, khăn ướt vắt quanh cổ, mồ hôi đã kịp khô. Vừa định khép mắt lại một chút để nghỉ ngơi thì Claire-san từ đâu bê ra một khay trà, mặt tươi như hoa.

``Xin lỗi đã để mọi người chờ! Mình đã chuẩn bị chút trà rồi đây,'' cô nàng nói, nhẹ nhàng như một bà sơ đúng nghĩa.

``Ừ, anh xin,'' tôi đáp, vươn tay cầm lấy một ly. ``À mà này\dots''

Tôi liếc quanh căn phòng. Dola thì đang ngồi gục trên sô-pha, miệng hé hé như sắp mớ một câu gì đó. Còn Ryushen-san thì tựa người bên kệ sách mới lắp, đầu gục xuống vai. Nghe kỹ, có thể nhận ra cả hai đang thở đều đều, thậm chí có tiếng khò khe khẽ. Tôi chỉ tay về phía họ.

\img{nj-4c}

``\dots À,'' Claire-san khẽ mỉm cười, hiểu ra.

``Chúng ta cũng sắp xong rồi,'' tôi nói, nhấp một ngụm trà ấm. ``Hay là nghỉ thôi nhỉ? Sáng mai ta tiếp tục.''

``Có vẻ là như vậy. Làm tốt lắm, mọi người,'' cô Nữ sơ ngồi xuống cạnh tôi, dịu dàng như mọi khi. Cô có vẻ đã quen với việc chăm sóc người khác đến mức tự nhiên như hơi thở.

Tôi gật đầu, uống cạn ly, rồi đặt nó xuống khay. ``OK\dots\ Vậy thì, anh xin phép.''

Tôi đứng dậy, bước ra cửa. Tay cầm lấy tay nắm cửa, kéo nhẹ. Bản lề cũ kỹ rên lên một tiếng khe khẽ, cứ như cũng hiểu được bầu không khí bình yên hiếm hoi này mà cố không làm phiền.

Cửa khép lại sau lưng tôi, và tôi rời khỏi văn phòng mới, nơi vừa bừa bộn vừa ấm áp như một ký ức mới hình thành.

Tôi không ở lại, nhưng vẫn có thể tưởng tượng ra cảnh tượng đằng sau cánh cửa ấy.

Claire-san - cô sơ nữ với mái tóc vàng và ánh mắt chan chứa lòng vị tha - chắc lại ngồi giữa hai đứa nhỏ đang ngủ kia. Có khi còn đặt đầu họ lên đùi mình, nhẹ nhàng vuốt tóc như một người mẹ ôm hai đứa con mệt mỏi. Có một thứ bình yên kỳ lạ toát ra từ cô ấy, kiểu yên ả khiến cả căn phòng lặng đi như thể thời gian ngừng trôi.

Phải rất lâu rồi tôi mới thấy được khoảnh khắc yên tĩnh như vậy. Ở \hll, ``yên lặng'' gần như là thứ xa xỉ. Ở đó luôn có tiếng hét của Fubu-chan, tiếng cười của Towa-chan, tiếng cãi nhau của Sui-chan và Miko-chan, tiếng nhạc mở sai beat của Nene-chan và đôi khi cả tiếng bàn phím đập vỡ từ A-san.

Nhưng ở đây, ít ra là lúc này, mọi thứ bỗng nhiên yên ắng đến lạ.

Tôi ngửa mặt lên trời, thở ra một hơi dài.

``\textit{Kỷ niệm\dots}'' tôi thì thầm, ``\textit{luôn là thứ khiến người ta muốn bước tiếp. Có lẽ mình cũng nên học cách quý nó hơn.}''
%==========================================TRUYỆN PHỤ=========================================
\omk

Tôi về đến nơi khi trời đã tối hẳn, ánh đèn đường hắt xuống con dốc phía trước cửa nhà tạo nên một thứ ánh sáng vàng vọt mà dễ chịu kỳ lạ. Vừa mở cửa, mùi thức ăn lập tức xộc vào mũi, một mùi thơm ấm cúng mà tôi đã quá quen trong những ngày vừa qua.

Trong bếp, Kaigou-chan đang xắn tay áo đứng cạnh bếp lò, còn Suzuki với cái tạp dề buộc gọn thì đang sắp bát đũa ra bàn. Hai cô gái ấy phối hợp với nhau như một cỗ máy trơn tru. Chắc chỉ cần mười phút là xong cả mâm cơm. Dù biết rằng hai chị em họ đã từng sống chung với nhau rất lâu nên việc phối hợp này khá là dễ hiểu, cái cảnh đó khiến tôi mỉm cười.

``\vie{Cậu về rồi à?}'' cô chị quay lại nhìn tôi. ``\vie{Cả ngày hôm nay cậu đi đâu thế?}''

``\vie{Ờ thì\dots\ phụ giúp bên phía mấy người bên Nijisanji chuyển đồ sang văn phòng mới,}'' tôi tháo giày và thở phào như thể vừa bỏ được một bao gạo khỏi vai.

``\vie{Thế sao anh không bảo tôi với Suzu đi cùng?}'' giọng cô có chút hờn nhẹ nhưng không trách móc, mà nghiêng về kiểu trách yêu thường thấy ở chị gái. ``\vie{Tôi bê nổi cả máy in hàng tấn đấy, còn Suzu thì khỏi nói, sắp đồ còn chuẩn hơn cả mấy phần mềm quản lý kho đó.}''

Tôi chỉ bật cười, lắc đầu nhẹ, rồi lười biếng đổ người xuống tấm nệm: ``\vie{Cũng không phải vấn đề gì to tát đâu. Chỉ là đôi lúc người ta cũng nên tự làm lấy vài việc, để mà giữ lại một chút gì đó gọi là\dots\ kỷ niệm.}''

Kaigou-chan chỉ giương mắt nhìn tôi với vẻ khó hiểu, rồi nhún vai quay lại chiếc nồi canh. Suzuki thì nhìn tôi, rồi nhẹ giọng: ``\vie{Kỷ niệm là gì ạ, Onii-sama?}''

Tôi chớp mắt, nhưng không trả lời ngay. Có lẽ tôi cũng chưa tìm được lời giải thích cho chính mình. Tôi chỉ đưa tay lên trán, che mắt, nhắm lại, thở đều và để mùi thức ăn dẫn lối tâm trí mình trôi lửng lơ.

Bữa tối hôm đó\dots\ không có tiếng nổ. Không có xiên oden phi thẳng vào nút khẩn cấp. Không có người bay lên trời hay khóc lóc vì một cái nơ cài tóc.

Chỉ có tiếng thìa va vào bát, tiếng nước súp nóng bốc hơi, và giọng Kaigou-chan nhắc Suzuki lấy thêm rau từ bếp phụ.

Lặng lẽ mà đầy đủ.

Sau ba ngày dồn dập như bị cuốn vào mớ bòng bong giữa các thế giới, hôm nay - lần đầu tiên - tôi mới thực sự cảm thấy mình có một buổi tối đúng nghĩa. Một buổi tối của việc được ``sống'', chứ không phải ``tồn tại.''

Và tôi ngủ sớm hơn thường lệ. Không mộng mị. Không giật mình giữa đêm. Không gì cả.

Chỉ đơn giản là một chút yên bình nhỏ nhoi hiếm có.

\end{document}